\subsection{Inequality of the mean}

We are going to sharpen Th. 1 for numerical measurable functions (finite or not).

DEFINITION 1. --- Let X be a locally compact space, \(\mu\) a measure on X. Given a numerical function f (finite or not), defined locally almost everywhere in X, we call maximum in measure, or \(\mu\) -maximum (resp. minimum in measure, or \(\mu\) -minimum) of the function f, and denote by \(\mathrm{M}_{\infty}(f)\) (resp. \(m_{\infty}(f)\) ), the infimum (resp. supremum) of the set of numbers \(\alpha\) such that \(f(x) \leqslant \alpha\) (resp. \(f(x) \geqslant \alpha\) ) locally almost everywhere (for \(\mu\) ).

It follows at once from the definition that \(m_{\infty}(f) = -M_{\infty}(-f)\) , thus from every property of the maximum in measure one deduces a corresponding property of the minimum in measure.

For every \(\alpha > \mathrm{M}_{\infty}(f)\) , the set of \(x \in \mathbf{X}\) such that \(f(x) > \alpha\) is locally negligible; now, the set of \(x \in \mathbf{X}\) such that \(f(x) > \mathrm{M}_{\infty}(f)\) is the union of the sets where \(f(x) > r_n\) , with \(r_n\) running over the set of rational numbers \(> \mathrm{M}_{\infty}(f)\) ; therefore \(f(x) \leqslant \mathrm{M}_{\infty}(f)\) locally almost everywhere (§5, No.~2). Similarly \(f(x) \geqslant m_{\infty}(f)\) locally almost everywhere; it follows that \(m_{\infty}(f) \leqslant \mathrm{M}_{\infty}(f)\) if the measure \(\mu\) is nonzero; moreover, the relation \(m_{\infty}(f) = \mathrm{M}_{\infty}(f)\) is equivalent to saying that \(f\) is equal to a constant locally almost everywhere. It is clear that if the measure \(\mu\) is nonzero, then

\[
\inf _ {x \in \mathrm{X}} f (x) \leqslant m _ {\infty} (f) \leqslant \mathrm{M} _ {\infty} (f) \leqslant \sup _ {x \in \mathrm{X}} f (x).
\]

If two functions \(f, g\) are equal locally almost everywhere, then \(m_{\infty}(f) = m_{\infty}(g)\) and \(\mathrm{M}_{\infty}(f) = \mathrm{M}_{\infty}(g)\) .

Finally, if \(f\) and \(g\) are two functions such that \(f + g\) is defined locally almost everywhere, then

\[
\mathrm{M} _ {\infty} (f + g) \leqslant \mathrm{M} _ {\infty} (f) + \mathrm{M} _ {\infty} (g)\tag{1}
\]

provided the second member is defined, as follows at once from Def. 1; similarly, if f and g are both \(\geqslant0\) and are such that fg is defined locally almost everywhere, then

\[
\mathrm{M} _ {\infty} (f g) \leqslant \mathrm{M} _ {\infty} (f) \mathrm{M} _ {\infty} (g)\tag{2}
\]

provided the second member is defined.

If \(M_{\infty}(f) < +\infty\) , then \(f(x) < +\infty\) locally almost everywhere, but not necessarily almost everywhere. A numerical function f is said to be bounded in measure (for the measure \(\mu\) ) if it is defined and finite almost everywhere and if, moreover, the numbers \(m_{\infty}(f)\) and \(M_{\infty}(f)\) are both finite (the latter condition amounts to saying that \(M_{\infty}(|f|)<+\infty\) ).

PROPOSITION 1 (Inequality of the mean). --- Let f be a measurable numerical function that is bounded in measure. For every integrable numerical function \(g \geqslant 0\) , the function fg (defined almost everywhere) is integrable and

\[
m _ {\infty} (f) \int g d | \mu | \leqslant \int f g d | \mu | \leqslant \mathrm{M} _ {\infty} (f) \int g d | \mu |.\tag{3}
\]

Moreover, two of the three members of the inequality (3) cannot be equal unless, in the set of \(x \in X\) such that \(g(x) \neq 0\) , \(f\) is equal to \(M_{\infty}(f)\) almost everywhere or equal to \(m_{\infty}(f)\) almost everywhere.

Indeed, \(fg\) is measurable (§5, No.~3, Cor. 5 of Th. 1); moreover, the inequality \(m_{\infty}(f)g(x) \leqslant f(x)g(x) \leqslant M_{\infty}(f)g(x)\) holds, not only locally almost everywhere, but even almost everywhere, because the set of points \(x \in X\) where \(g(x) \neq 0\) is a countable union of integrable sets (§5, No.~6, Lemma 1). It follows that \(fg\) is integrable (§5, No.~6, Th. 5) and the inequality (3) holds. On the other hand, the function \(M_{\infty}(f)g - fg\) is almost everywhere defined and equal to \((M_{\infty}(f) - f)g\) ; it is therefore \(\geqslant 0\) almost everywhere in \(X\) ; since the relation \(M_{\infty}(f) \int g d|\mu| = \int fg d|\mu|\) is equivalent to \(\int(M_{\infty}(f) - f)gd|\mu| = 0\) , it can hold only if the function \((M_{\infty}(f) - f)g\) is negligible, which completes the proof.

Setting aside the trivial case that \(\int g d|\mu| = 0\) , the inequality (3) may be deduced from Th. 1 of No.~1 applied to the interval \(D = [m_{\infty}(f), M_{\infty}(f)]\) . One can bring to Th. 1 of No.~1 the complements analogous to those of Prop. 1, that specify the case in which the point \((\int f g d\mu)/(\int g d\mu)\) belongs to the boundary of D (Exer. 2).

